% Generated draft: requires mathematical and source review.
\documentclass[11pt]{article}
\usepackage[margin=25mm]{geometry}
\usepackage{fontspec}
\setmainfont{Latin Modern Roman}
\usepackage{amsmath,amssymb,mathtools,unicode-math}
\setmathfont{Latin Modern Math}
\usepackage{enumitem,needspace,xurl,hyperref}
\setlength{\parindent}{0pt}
\setlength{\parskip}{5pt}
\setlength{\emergencystretch}{3em}
\newcommand{\N}{\mathbb{N}}
\newcommand{\Z}{\mathbb{Z}}
\newcommand{\Q}{\mathbb{Q}}
\newcommand{\R}{\mathbb{R}}
\newcommand{\C}{\mathbb{C}}
\newcommand{\F}{\mathbb{F}}
\newcommand{\A}{\mathbb{A}}
\newcommand{\PP}{\mathbb{P}}
\newcommand{\E}{\mathbb{E}}
\newcommand{\eps}{\varepsilon}
\newcommand{\dd}{\,\mathrm{d}}
\begin{document}
\section*{UnsolvedMath ID: 20003092}
\textbf{Source problem code:} catalogue entry
\subsection{Definitions and mathematical statement}
Let $S$ be an orientable finite-type surface of negative Euler characteristic in the intended scope, and let $\mathcal T(S)$ be its Teichm\"uller space. The directed Thurston metric is
\[
 d_{\mathrm{Th}}(X,Y)=\log\sup_{\alpha}
       \frac{\ell_\alpha(Y)}{\ell_\alpha(X)},
\]
where $\alpha$ ranges over essential simple closed curves. A directed geodesic segment $\gamma:[0,D]\to\mathcal T(S)$ from $X$ to $Y$ satisfies $D=d_{\mathrm{Th}}(X,Y)$ and $d_{\mathrm{Th}}(\gamma(s),\gamma(t))=t-s$ for $s\leq t$. Problem 2.2 asks to describe
\[
 \operatorname{Env}(X,Y)=\bigcup_{\gamma:X\to Y}\gamma([0,D]).
\]
The union is directed from $X$ to $Y$; it need not equal the reverse envelope. The source does not restrict $S$ to a particular complexity or closedness class.

The supplied research assessment gives the following scopes. Dumas--Lenzhen--Rafi--Tao describe once-punctured-torus envelopes as either a unique stretch segment or a compact geodesic quadrilateral. Bar-Natan's 2023 version-1 preprint gives an analogous unique-geodesic/combinatorial-quadrilateral description for the four-punctured sphere. Pan--Wolf, \texttt{arXiv:2401.06607v2}, revised January 21, 2026, describe envelopes for closed orientable surfaces of genus at least two as contractible iterated cones, a cone over a cone over an ``ancestor'' determined by strict chain-recurrent extensions of the maximally stretched lamination. The summary does not define this ancestor or the precise cone parameterization, so those terms require the cited theorem. It explicitly says the Pan--Wolf methods do not cover punctured surfaces and that no checked source in that assessment supplies the shape theorem for all higher-complexity punctured types.

The aid adds a general directed-interval structure. Let $(M,d)$ be a geodesic asymmetric metric space with $d\geq0$ and $d(P,Q)=0$ only for $P=Q$. Then
\[
 \operatorname{Env}(X,Y)=
 \{P:d(X,Y)=d(X,P)+d(P,Y)\}.
\]
On this set put $h(P)=d(X,P)$ and define
\[
 P\preceq Q\quad\Longleftrightarrow\quad
 h(Q)=h(P)+d(P,Q).
\]
The aid states that this is a partial order graded by $h$, with
\[
 \{P:P\preceq Q\}=\operatorname{Env}(X,Q),\qquad
 \{Q:P\preceq Q\}=\operatorname{Env}(P,Y).
\]
For comparable $P\preceq Q$, its order interval is exactly $\operatorname{Env}(P,Q)$. Every directed $X$-to-$Y$ geodesic image is a maximal chain whose rank image is the full interval $[0,d(X,Y)]$. This does not assert that an arbitrary maximal chain has full rank without additional hypotheses. These metric-order statements organize an envelope but do not determine its surface-specific topology or settle the remaining punctured-surface problem.
\subsection{Short English statement}
Describe the union of all directed Thurston geodesics between two marked surfaces. The supplied assessment gives explicit low-complexity descriptions and a January 2026 closed-surface shape theorem, with punctured-surface limitations. A separate general reduction orders an envelope by distance from its initial point and identifies its subintervals with smaller envelopes.
\subsection{Sources}
\begin{itemize}
\item Local UnsolvedMath registry, permanent ID $20003092$, source code catalogue entry; canonical location \texttt{aim-topology-notes.json}, \texttt{notes[179]}.
\item Weixu Su, AIM list \emph{Problems on Thurston Metric}, April 21, 2014, Problem 2.2, \url{https://aimath.org/pastworkshops/lipschitzteichproblems.pdf}. The extraction attributes it to ``F. Gu\'eritau''; its exact spelling is unresolved here. The supplied verification reports direct PDF inspection, not performed during this draft.
\item Dumas--Lenzhen--Rafi--Tao; Bar-Natan, 2023 version-1 preprint; and Pan--Wolf, \texttt{arXiv:2401.06607v2}, revised January 21, 2026, as named in the supplied research summary. No titles or further identifiers are supplied for the first two references.
\item Research attempt 1 and summary are marked partial result, proved, and low novelty confidence, with an explicit machine-generated, non-peer-reviewed research-aid caveat requiring independent verification before citation.
\end{itemize}
\end{document}
